\section*{अध्याय \ref{ch:b3:lp} --- $L^p$ समष्टियाँ}

\begin{solution}{exo:b3:lp:1}
(क) \cref{thm:b3:lp:holder} में $p = q = 2$: $\abs{\int f\bar
g\,\dd\mu} \leq \int\abs{fg} \leq \norm f_2\norm g_2$ --- अर्थात् कोशी--श्वार्ज़, जिसमें बराबरी तभी
और केवल तभी जब $\abs f, \abs g$ समानुपाती हों और कलाएँ संरेखित।

(ख) किसी प्रायिकता समष्टि पर $p \leq q$ के लिए: उत्तल $\Phi(t) =
\abs t^{q/p}$ के साथ येंसन
(\cref{exo:b3:lp:10}) को फलन $\abs f^p$ पर लगाइए: $\bigl(\int\abs f^p\bigr)^{q/p} \leq \int\abs f^q$, अर्थात् $\norm f_p \leq \norm f_q$।

(ग) $\int\abs{fg} = \norm f_1\norm g_\infty$ तभी और केवल तभी जब $\{f \neq 0\}$ पर लगभग सर्वत्र $\abs{g} =
\norm g_\infty$ (असमिका $\abs{fg} \leq \abs f\norm g_\infty$
को लगभग सर्वत्र बराबरी होना चाहिए)।
\end{solution}

\begin{solution}{exo:b3:lp:2}
$\int_0^1 x^{-p/2}\dd x < \infty$ तभी और केवल तभी जब $p < 2$: अतः पहला $p \in \intco12$ के लिए $L^p$ में है।
$\int_1^\infty x^{-p/2}\dd x <
\infty$ तभी और केवल तभी जब $p > 2$: अतः दूसरा $p \in \intoo2\infty$ के लिए (और $p = \infty$: वह
परिबद्ध है --- उसे सम्मिलित कीजिए)। तीसरा: $p <
2$ के लिए $x^{-p/2}$ का
प्रभुत्व: \omterm{def:b3:lebesgue:l1}{समाकलनीय}; $p = 2$ के लिए $u = -\ln x$ प्रतिस्थापित कीजिए: $\int_0^1\frac{\dd x}{x(1 + \abs{\ln
x})^2} = \int_0^\infty\frac{\dd u}{(1 + u)^2} < \infty$; और
$p > 2$ के लिए घात हावी हो जाती है: अपसारी। अतः $p \in \intcc12$। चौथा: $\int_\R\frac{\dd x}{(1 + \abs x)^p} < \infty$ तभी
और केवल तभी जब $p >
1$; परिबद्ध, अतः $L^\infty$ भी: $p \in
\intoc1\infty$। शिक्षा: $0$ पर
\omterm{def:b3:lebesgue:l1}{समाकलनीयता} को छोटे $p$ भाते हैं, और $\infty$ पर बड़े $p$; दोनों बाधाओं
को मिलाने पर $L^p(\R)$ समष्टियों के बीच कोई अंतर्विष्टि नहीं।
\end{solution}

\begin{solution}{exo:b3:lp:3}
(क) $\norm{f_n}_p^p = \lambda(I_n) \to 0$ (द्विअंशी स्तर $k$ पर लंबाई $2^{-k}$ है)। पर प्रत्येक $x$
प्रत्येक द्विअंशी स्तर के किसी एक अंतराल में होता है: अतः अनंत बार
$f_n(x) = 1$ और अनंत बार $= 0$ ($x$ को न धारण करने वाले उसी स्तर के अंतराल):
इसलिए किसी भी बिंदु पर अभिसरण नहीं।

(ख) $f_{n_k} = \mathbf 1_{\intcc0{2^{-k}}}$ (प्रत्येक स्तर का पहला अंतराल) प्रत्येक $x > 0$ पर $0$ पर
अभिसरित होता है: अर्थात् लगभग सर्वत्र।

(ग) लगभग सर्वत्र पर $L^1$ में नहीं: लगभग सर्वत्र $n\mathbf 1_{\intoo0{1/n}} \to 0$, समाकल $1$।
$L^1$ में पर किसी $L^p$ में नहीं ($p > 1$): $g_n =
\eu^n\,\mathbf 1_{(0,\ \eu^{-n}/n)}$: $\norm{g_n}_1 = \frac1n
\to 0$, जबकि प्रत्येक
$p > 1$ के लिए $\norm{g_n}_p^p = \eu^{(p-1)n}/n \to \infty$।
\end{solution}

\begin{solution}{exo:b3:lp:4}
ऊपरी: $\norm f_p \leq \mu(X)^{1/p}\norm f_\infty$ ($q = \infty$ के साथ \cref{prop:b3:lp:inclusions}(क)), और $\mu(X)^{1/p} \to 1$। निचली: $\varepsilon > 0$ के लिए $A =
\{\abs f > \norm f_\infty - \varepsilon\}$
के पास $\mu(A) > 0$ है (सारभूत उच्चतम की परिभाषा), और
\[
\norm f_p \geq \Bigl(\int_A \abs f^p\Bigr)^{1/p}
\geq (\norm f_\infty - \varepsilon)\,\mu(A)^{1/p}
\xrightarrow[p\to\infty]{} \norm f_\infty - \varepsilon .
\]
\end{solution}

\begin{solution}{exo:b3:lp:5}
(क) दो बार: रूपांतरण $\norm{\tau_hg - g}_p \leq
C^{1/p}\norm{\tau_hg - g}_\infty$ (परिमित-माप आलंब) $p = \infty$ के लिए केवल इस
अर्थ में \omterm{rem:b3:complete:meagre}{क्षीण} हो जाता है कि वहाँ \emph{$\mathcal C_c$ की सघनता} विफल हो जाती
है --- और वही असली दरार है: \cref{thm:b3:lp:density} के चरण (2) का कोई $L^\infty$ सादृश्य
नहीं है।

(ख) यदि $f$ का कोई एकसमान \omterm{def:b3:topology:continuity}{संतत} प्रतिनिधि $g$ हो: $\norm{\tau_hf - f}_\infty = \sup_x\abs{g(x - h) - g(x)} \to
0$। विलोमतः
मान लीजिए $\norm{\tau_hf - f}_\infty \to 0$। मृदुकरण $f_\varepsilon = f * \rho_\varepsilon$ \omterm{def:b3:topology:continuity}{संतत} हैं, और
\[
\norm{f_\varepsilon - f}_\infty
\leq \sup_{\norm y \leq \varepsilon}\norm{\tau_yf -
f}_\infty \longrightarrow 0
\]
(संवलन स्थानांतरों का कोई औसत है)। प्रत्येक $f_\varepsilon$ एकसमान रूप से \omterm{def:b3:topology:continuity}{संतत}
है (औसत लेने से $\norm{\tau_hf_\varepsilon - f_\varepsilon}_\infty \leq
\norm{\tau_hf - f}_\infty$), और एकसमान \omterm{def:b3:topology:continuity}{संतत} फलनों की कोई एकसमान सीमा भी
वैसी ही होती है: अतः $f$ किसी एकसमान \omterm{def:b3:topology:continuity}{संतत} फलन से लगभग सर्वत्र सहमत
है।
\end{solution}

\begin{solution}{exo:b3:lp:6}
होल्डर से प्रत्येक $x$ के लिए समाकल्य $y \mapsto f(x -
y)g(y)$ $\abs{(f*g)(x)} \leq \norm f_p\norm
g_q$ के साथ $L^1$ में
है: अतः सर्वत्र परिभाषित और परिबद्ध। एकसमान \omterm{def:b3:topology:continuity}{सांतत्य} ($p <
\infty$):
\[
\abs{(f*g)(x + h) - (f*g)(x)}
= \Bigl|\int\bigl(\tau_{-h}f - f\bigr)(x - y)\,g(y)\dd y\Bigr|
\leq \norm{\tau_{-h}f - f}_p\,\norm g_q
\xrightarrow[h\to0]{} 0,
\]
$x$ में एकसमान रूप से (\cref{thm:b3:lp:density}(3))। यदि $p =
\infty$ हो, तो $q = 1$: $f * g = g * f$ लिखिए
और वही परिबंध चलाइए, जिसमें स्थानांतरण $g \in L^1$ पर क्रिया करता है।
\end{solution}

\begin{solution}{exo:b3:lp:7}
$\int\chi =
1$ वाला $\chi \in \mathcal C_c^\infty(\intoo ab)$ स्थिर कीजिए, और $c = \int f\chi$ रखिए। मान लीजिए $\varphi \in \mathcal
C_c^\infty$ स्वेच्छ
है और $\psi = \varphi -
\bigl(\int\varphi\bigr)\chi$: तब $\int\psi = 0$, अतः $\Phi(x)
= \int_a^x\psi$ $\Phi' = \psi$ के साथ $\Phi \in \mathcal C_c^\infty(\intoo
ab)$ परिभाषित करता
है (वह दोनों सिरों के निकट लुप्त हो जाता है: $a$ के निकट तुच्छ
रूप से, और $b$ के निकट इसलिए कि कुल समाकल $0$ है)। परिकल्पना
$\int f\psi = \int f\Phi' = 0$ देती है, अतः
\[
\int f\varphi = \Bigl(\int\varphi\Bigr)\int f\chi = \int
c\,\varphi
\quad\text{प्रत्येक } \varphi \text{ के लिए}:
\qquad \int(f - c)\varphi = 0 .
\]
\cref{cor:b3:lp:fundlemma} ($\intoo ab$ पर स्थानीयकृत) से लगभग सर्वत्र $f =
c$।
\end{solution}

\begin{solution}{exo:b3:lp:8}
मान लीजिए $3\delta < d(K, \R^d\setminus U)$ (धनात्मक: \cref{exo:b3:topology:6}(ख); यदि $U = \R^d$ हो तो कोई भी $\delta$
काम करता है), $K_\delta = \{x : d(x, K) \leq \delta\}$, और
\[
\varphi = \mathbf 1_{K_\delta} * \rho_{\delta/2} .
\]
तब $\varphi \in \mathcal C^\infty$ (\cref{thm:b3:lp:regularization}(1); सूचक $L^1$ है), $0 \leq \varphi \leq 1$ ($\int\rho = 1$), $K$ पर $\varphi = 1$
($x \in K$ के लिए $\bar B(x, \delta/2) \subseteq K_\delta$, अतः संवलन $\rho$ का पूरा समाकलन कर लेता है), और
$\operatorname{supp}\varphi \subseteq K_{3\delta/2} \subseteq
U$: अर्थात् \omterm{def:b3:topology:compact}{संहत} आलंबित ($K_\delta$ परिबद्ध है)। इकाई विभाजन: $K \subseteq U_1\cup\dots\cup U_m$
दिया हो, तो (संहतता से) $K \subseteq
\bigcup \mathring K_i$ वाले \omterm{def:b3:topology:compact}{संहत} $K_i \subseteq U_i$ चुनिए, $(K_i,
U_i)$ के लिए
ऊपर की भाँति $\varphi_i$ लीजिए, और $\psi_i = \varphi_i\prod_{j <i}(1 -
\varphi_j)$ रखिए: प्रत्येक $\psi_i \in \mathcal C_c^\infty(U_i)$, और $K$
पर $\sum_i\psi_i = 1 - \prod_i(1 - \varphi_i) = 1$।
\end{solution}

\begin{solution}{exo:b3:lp:9}
(क) $r = p$ पर $(p_0, q_0) = (1, \infty)$ के लिए अंतर्वेशन घातांक $\alpha = \frac1p$ है: \cref{prop:b3:lp:inclusions}(ग) $\norm f_p \leq \norm
f_1^{1/p}\norm f_\infty^{1 - 1/p}$
दे देता है।

(ख) \cref{prop:b3:lp:inclusions}(ग) में लघुगणक लेने पर: $\ln\norm f_r \leq
\alpha\ln\norm f_p + (1-\alpha)\ln\norm f_q$, जहाँ $\frac1r$ $\frac1p, \frac1q$ का वही
उत्तल संचय है: अतः $\frac1p
\mapsto \ln\norm f_p$ उत्तल है। ठीक $\intoo{p_0}{p_1}$ पर $L^p$-सदस्यता वाला
उदाहरण:
\[
f(x) = x^{-1/p_1}\,\mathbf 1_{\intoo01}(x) +
x^{-1/p_0}\,\mathbf 1_{\intco1\infty}(x):
\]
पहला पद $L^p$ तभी और केवल तभी जब $p < p_1$, और दूसरा तभी जब $p >
p_0$।
\end{solution}

\begin{solution}{exo:b3:lp:10}
मान लीजिए $m = \int f\,\dd\mu \in \R$। उत्तलता $m$ पर कोई आधार रेखा देती है: ऐसा $s$
है कि सभी $t$ के लिए $\Phi(t) \geq \Phi(m) + s(t -
m)$ ($s$ को एकपक्षीय अवकलजों के बीच लीजिए,
जो उत्तल फलनों के लिए विद्यमान हैं)। $t = f(x)$ प्रतिस्थापित कीजिए और
प्रायिकता $\mu$ के सापेक्ष समाकलन कीजिए:
\[
\int\Phi\circ f\,\dd\mu \geq \Phi(m) + s\Bigl(\int f - m\Bigr)
= \Phi\Bigl(\int f\,\dd\mu\Bigr)
\]
(\omterm{def:b3:lebesgue:measurable}{मापनीयता}: $\Phi$ \omterm{def:b3:topology:continuity}{संतत} है; और $\Phi\circ f$ के ऋणात्मक भाग की \omterm{def:b3:lebesgue:l1}{समाकलनीयता} आधार
रेखा से प्रत्याभूत है)। समांतर माध्य--गुणोत्तर माध्य: भार $w_i$ वाले
किसी परिमित समुच्चय पर $\Phi = \exp$ और $f = \sum(\ln a_i)\mathbf 1_i$ लीजिए: $\exp\bigl(\sum w_i\ln a_i\bigr) \leq \sum w_ia_i$, अर्थात् $\prod a_i^{w_i} \leq \sum w_ia_i$।
मानक की एकदिष्टता \cref{exo:b3:lp:1}(ख) है।
\end{solution}

\begin{solution}{exo:b3:lp:11}
(क) पहले $p, q, r < \infty$ और $f, g \geq 0$ मान लीजिए (निरपेक्ष मानों से बदल दीजिए)।
तीनों घातांक $r$, $\alpha = \frac{pr}{r-p}$, $\beta = \frac{qr}{r-q}$ $\frac1r + \frac1\alpha + \frac1\beta = \frac1r +
\frac1p - \frac1r + \frac1q - \frac1r = 1$ संतुष्ट करते हैं (मापन संबंध)।
स्थिर $x$ के लिए विभाजित कीजिए,
\[
f(y)g(x-y) = \bigl[f(y)^pg(x-y)^q\bigr]^{1/r}\cdot
f(y)^{1 - p/r}\cdot g(x-y)^{1 - q/r},
\]
और तीनों घातांकों के साथ होल्डर
\[
(f*g)(x) \leq \Bigl(\int f^pg(x-\cdot)^q\Bigr)^{1/r}
\norm f_p^{\,p/\alpha\cdot\alpha/p}\cdots
\]
दे देता है; अधिक परिशुद्धता से: दूसरा गुणक $\bigl(\int f^{(1-p/r)\alpha}\bigr)^{1/\alpha} =
\norm f_p^{p(1/p - 1/r)}$ है क्योंकि $(1 - \frac pr)\alpha = p$,
और इसी प्रकार तीसरा $\norm g_q^{q(1/q - 1/r)}$ है। $r$-वीं घात लीजिए और $x$ में समाकलन
कीजिए (पहले गुणक पर टोनेली):
\[
\norm{f*g}_r^r \leq \norm f_p^p\,\norm g_q^q\cdot
\norm f_p^{\,rp(1/p - 1/r)}\,\norm g_q^{\,rq(1/q - 1/r)}
= \norm f_p^r\,\norm g_q^r .
\]
अंत-बिंदु स्थितियाँ ($r = \infty$ अथवा कोई घातांक अपनी सीमा के बराबर) सादा
होल्डर या सीधे आकलन हैं।

(ख) $r = \infty$ $q = p'$ अनिवार्य कर देता है: $\abs{f*g(x)} \leq \norm
f_p\norm g_{p'}$ --- अर्थात् स्थानांतरण-परावर्तन
के बाद होल्डर। $q = 1$ $r = p$ देता है: $\norm{f*g}_p \leq \norm
g_1\norm f_p$, अर्थात् मृदुकरण का भारवाहक
(\cref{thm:b3:lp:regularization} का इंजन)। $p = q = 1$ $r = 1$ देता है: अर्थात् संवलन बीजगणित (\cref{thm:b3:product:convolution})।

(ग) $f, g$ को $f_\lambda = f(\lambda\cdot)$ से और $g_\lambda = g(\lambda\cdot)$ को बदल दीजिए: तब $f_\lambda*g_\lambda =
\lambda^{-d}(f*g)(\lambda\cdot)$, और मानकों
की तुलना करने पर
\[
\text{बायाँ पक्ष} \sim \lambda^{-d - d/r}, \qquad
\text{दायाँ पक्ष} \sim \lambda^{-d/p - d/q} :
\]
सभी $f, g$ के लिए वैध कोई असमिका दोनों मापन घातांकों को मेल खाने पर
बाध्य कर देती है, अर्थात् ठीक $1 + \frac1r = \frac1p + \frac1q$। कोई भी अन्य संयोजन $\lambda \to 0$ या
$\infty$ पर मर जाता है: यंग का संबंध कोई सुविधा नहीं, बल्कि एक मापन
नियम है।
\end{solution}

\begin{solution}{exo:b3:lp:12}
(क) $\norm f_p = \norm g_q = 1$ मानकीकृत कीजिए। होल्डर की उपपत्ति यंग की असमिका $\abs{fg} \leq
\frac{\abs f^p}p + \frac{\abs g^q}q$ का
समाकलन करती है; और समाकलों की बराबरी यंग में लगभग सर्वत्र बराबरी
अनिवार्य कर देती है, जिसका ($\exp$ की कड़ी उत्तलता; बराबरी तभी और
केवल तभी जब $a^p = b^q$) अर्थ है लगभग सर्वत्र $\abs
f^p = \abs g^q$। मानकीकरण हटाने पर:
लगभग सर्वत्र $\abs f^p\norm g_q^q = \abs g^q\norm f_p^p$ --- अर्थात् समानुपातिकता।

(ख) मिन्कोव्स्की $\abs{f + g}^{p-1}
\abs f$ और $\abs{f+g}^{p-1}\abs g$ पर लगाए गए दो होल्डर हैं; और बराबरी
दोनों में (क) वाली समानुपातिकता अनिवार्य कर देती है: $\abs f^p$ और $\abs g^p$
प्रत्येक $\abs{f+g}^{(p-1)q} = \abs{f+g}^p$ के समानुपाती, अतः किसी अचर $t \geq 0$ के लिए लगभग सर्वत्र
$\abs g = t\abs f$; और प्रारंभिक बिंदुशः त्रिभुज असमिका $\abs{f + g} \leq
\abs f + \abs g$ को भी लगभग सर्वत्र
बराबरी होना चाहिए, जिसका सम्मिश्र मानों के लिए अर्थ है कि जहाँ दोनों
शून्येतर हों वहाँ $f$ और $g$ का कोणांक लगभग सर्वत्र एक ही हो।
मिलाकर: लगभग सर्वत्र $g = tf$, $t > 0$ (दोनों शून्येतर)।

(ग) $p = 1$: $\int\abs{f + g} = \int\abs f +
\int\abs g$ में बराबरी तब भी लागू रहती है जब $f, g$ का चिह्न
प्रतिरूप एक ही हो (लगभग सर्वत्र एक ही कोणांक) --- किसी समानुपातिकता
की आवश्यकता नहीं: $f = \mathbf 1_{\intcc01}$ और $g = \mathbf 1_{\intcc02}$ काम कर देते हैं। $p = \infty$: $\norm{f+g}_\infty = \norm f_\infty +
\norm g_\infty$ तभी
हो जाता है जब दोनों फलन किसी उभयनिष्ठ बिंदु पर (अथवा किसी उभयनिष्ठ
अनुक्रम के अनुदिश) संगत रूप से शिखर बनाएँ: अन्यत्र के व्यवहार की
परवाह किए बिना किसी एक बिंदु के निकट $f = g\,$ पर्याप्त है। $1 <
p < \infty$ के
लिए $L^p$ गोलों की कड़ी उत्तलता --- और अंत-बिंदुओं पर उसकी विफलता
--- ठीक वही है जिसकी ये बराबरी स्थितियाँ साक्षी हैं।
\end{solution}

\begin{solution}{pb:b3:lp:1}
\textbf{1.} $f$ का आलंब किसी $[\alpha, \beta]
\subseteq \intoo0\infty$ में है, अतः $[0, \alpha]$ पर $F = 0$ और
$[\beta, \infty)$ पर $F
\equiv F(\beta)$: इसलिए $Hf$ $0$ के निकट लुप्त हो जाता है और अनंत
पर $O(1/x)$ है; $p > 1$ के लिए $\int_1^\infty x^{-p}\dd x <
\infty$, और $Hf$ \omterm{def:b3:topology:continuity}{संतत} है: $Hf \in L^p$।

\textbf{2.} $\bigl(x^{1-p}F(x)^p\bigr)' = (1-p)x^{-p}F^p +
p\,x^{1-p}F^{p-1}f$। दोनों परिसीमा मान लुप्त हो जाते हैं: $0$ पर
इसलिए कि $0$ के निकट $F = 0$; और $\infty$ पर इसलिए कि $x^{1-p}F^p \leq
F(\beta)^p x^{1-p} \to 0$ ($p > 1$)।
इस सर्वसमिका का $\intoo0\infty$ पर समाकलन:
\[
0 = (1 - p)\int_0^\infty\Bigl(\frac Fx\Bigr)^p\dd x
+ p\int_0^\infty\Bigl(\frac Fx\Bigr)^{p-1}f(x)\,\dd x,
\]
जो दिखाया गया संबंध है।

\textbf{3.} घातांकों $q = \frac p{p-1}$ और $p$ के साथ होल्डर:
\[
\int\Bigl(\frac Fx\Bigr)^{p-1}f
\leq \Bigl(\int\Bigl(\frac
Fx\Bigr)^{p}\Bigr)^{1 - 1/p}\,\norm f_p ,
\]
अतः $\norm{Hf}_p^p \leq \frac p{p-1}\norm{Hf}_p^{p-1}\norm
f_p$; भाग दीजिए (प्रश्न 1 से परिमित, और यदि $0$ हो तो सिद्ध
करने को कुछ नहीं)।

\textbf{4.} मान लीजिए $f \in L^p$, $f \geq 0$। $L^p$ में $\varphi_k \to f$ वाला $\varphi_k \in
\mathcal C_c(\intoo0\infty)$ चुनिए
(\cref{thm:b3:lp:density}(2), विवृत अर्ध-रेखा से प्रतिच्छेदित --- $f\mathbf 1_{[1/k, k]}$ का आसन्नन कीजिए
और विकर्ण लीजिए), और $\varphi_k$ को $\abs{\varphi_k}$ से बदल दीजिए (अब भी \omterm{def:b3:topology:continuity}{संतत}, और
$f \geq 0$ के अधिक निकट)। प्रत्येक स्थिर $x > 0$ के लिए $\intoo0x$ पर होल्डर
\[
\abs{H\varphi_k(x) - Hf(x)} \leq
\frac1x\,x^{1 - 1/p}\,\norm{\varphi_k - f}_p \to 0 :
\]
देता है, और बिंदुशः $H\varphi_k \to Hf$। फतू और प्रश्न 3:
\[
\int (Hf)^p \leq \liminf_k\int(H\varphi_k)^p
\leq \Bigl(\frac p{p-1}\Bigr)^p\liminf_k\norm{\varphi_k}_p^p
= \Bigl(\frac p{p-1}\Bigr)^p\norm f_p^p .
\]
चिह्नयुक्त या सम्मिश्र $f$ के लिए: बिंदुशः $\abs{Hf} \leq H\abs f$, और ऋणेतर स्थिति
$\abs f$ पर लागू होती है।

\textbf{5.} $\norm{f_A}_p^p = \int_1^A t^{-1}\dd t = \ln A$। $1 \leq x \leq A$ के लिए:
\[
(Hf_A)(x) = \frac1x\int_1^x t^{-1/p}\dd t
= \frac{x^{1 - 1/p} - 1}{(1 - \tfrac1p)\,x}
= \frac p{p-1}\,x^{-1/p}\bigl(1 - x^{-(1 - 1/p)}\bigr).
\]

\textbf{6.} $\varepsilon > 0$ और $x \geq X_0$ के लिए $(1 -
x^{-(1-1/p)})^p \geq 1 - \varepsilon$ वाला $X_0$ स्थिर कीजिए। तब
\[
\norm{Hf_A}_p^p \geq \int_{X_0}^A\Bigl(\frac
p{p-1}\Bigr)^p\frac{1 - \varepsilon}{x}\,\dd x
= \Bigl(\frac p{p-1}\Bigr)^p(1 - \varepsilon)\,(\ln A - \ln
X_0),
\]
अतः $A \to \infty$ पर $\dfrac{\norm{Hf_A}_p^p}{\norm{f_A}_p^p} \geq \bigl(\frac
p{p-1}\bigr)^p(1 - \varepsilon)\bigl(1 - \frac{\ln X_0}{\ln
A}\bigr) \to \bigl(\frac p{p-1}\bigr)^p(1 - \varepsilon)$: इसलिए अचर को सुधारा नहीं जा सकता।

\textbf{7.} प्रश्न 3 में बराबरी होल्डर में बराबरी अनिवार्य कर देती
है: अर्थात् लगभग सर्वत्र $f^p$ $\bigl(\frac Fx\bigr)^{(p-1)q}
= \bigl(\frac Fx\bigr)^p$ का समानुपाती, यानी किसी $\gamma \geq 0$
के लिए लगभग सर्वत्र $f = \gamma\,\frac Fx$। चूँकि $F(x) = \int_0^xf$ निरपेक्षतः \omterm{def:b3:topology:continuity}{संतत} है और लगभग
सर्वत्र $F' = f$, अतः $F$ $F' =
\gamma F/x$ हल करता है: जिस भी अंतराल पर $F > 0$
हो वहाँ $(\ln F)' =
\gamma/x$, अतः $F = c\,x^{\gamma}$ और वहाँ $f = c\gamma
x^{\gamma - 1}$। पर कोई शून्येतर घात $x^{\gamma-1}$
$L^p(\intoo0\infty)$ में नहीं होती ($\infty$ पर $p(\gamma - 1) < -1$ चाहिए और $0$ पर $> -1$: असंगत),
और $F$ सर्वथा लुप्त नहीं हो सकता जब तक $f = 0$ न हो। अतः बराबरी
$f =
0$ माँगती है।

\textbf{8.} अनवर्धमान $\geq 0$ $(a_n)$ दिया हो, तो $\intoo0{+\infty}$ पर सोपान फलन
$g(t) = a_{\lceil t\rceil}$ परिभाषित कीजिए: यह अनवर्धमान है, जिसमें $\int_0^\infty g^p =
\sum_na_n^p$ और $\int_0^ng = a_1 + \dots + a_n$, अतः
$(Hg)(n) = \frac{a_1 + \dots + a_n}n$। किसी अनवर्धमान फलन का औसत अनवर्धमान होता है, अतः $Hg$ भी
वैसा ही है, और भाग I से
\[
\sum_{n\geq1}\Bigl(\frac{a_1{+}\dots{+}a_n}n\Bigr)^p
= \sum_{n\geq1}(Hg)(n)^p
\leq \sum_{n\geq1}\int_{n-1}^n (Hg)(t)^p\,\dd t
= \norm{Hg}_p^p
\leq \Bigl(\frac p{p-1}\Bigr)^p\sum_n a_n^p
\]
किसी व्यापक ऋणेतर अनुक्रम के लिए मान लीजिए $(a_n^*)$ उसका अनवर्धमान
पुनर्विन्यास है (यह $a_n \to 0$ होने पर संभव है, और हम वैसा मान सकते हैं
--- अन्यथा दोनों पक्ष अनंत हैं): दायाँ पक्ष अपरिवर्तित रहता है, और
प्रत्येक आंशिक योग $a_1 + \dots
+ a_n$ अधिक से अधिक $a_1^* + \dots + a_n^*$ है ($n$ सबसे बड़े पद):
अतः बायाँ पक्ष केवल बढ़ता है। इसलिए सभी $(a_n)$ के लिए असमिका।

\textbf{9.} यदि $(a_n) \in \ell^p$ हो, तो चेज़ारो माध्यों का अनुक्रम मानक $\leq \frac p{p-1}\norm
a_p$
के साथ $\ell^p$ में है। उदाहरण ($p = 2$): $a_n = \frac1{\sqrt n\,\ln n}$ ($n
\geq 2$): $\sum a_n^2 = \sum\frac1{n\ln^2n} < \infty$, फिर भी
$\sum a_n = \infty$ (समाकल परीक्षा); हार्डी फिर भी $\sum_n\bigl(\frac{a_1 + \dots + a_n}n\bigr)^2 < \infty$ की प्रत्याभूति देते हैं।

\textbf{10.} $f = \mathbf 1_{\intcc01}$ के लिए: $\intoc01$ पर $Hf(x) = 1$ और $x \geq 1$ के लिए $= \frac1x$:
$\int_0^\infty Hf =
1 + \int_1^\infty\frac{\dd x}x = \infty$, जबकि $\norm f_1 =
1$। उपपत्ति दो स्थानों पर ढह जाती है: अचर $\frac p{p-1}$ $p \to 1$
पर फट जाता है, और परिसीमा पद $x^{1-p}F^p$ $p = 1$ के लिए अनंत पर अब लुप्त
नहीं होता। हार्डी की असमिका एक ईमानदार $p > 1$ परिघटना है।

\textbf{11.} सबसे लंबा अंतराल $B_{i_1}$ चुनिए; उससे मिलने वाला प्रत्येक
अंतराल हटा दीजिए; बचे हुओं में से सबसे लंबा $B_{i_2}$ चुनिए; और पुनरावृत्ति
कीजिए (अंतराल परिमित संख्या में हैं)। चुने गए रचना से असंयुक्त हैं,
और हटाया गया प्रत्येक $B$ कम से कम उतने ही लंबे किसी चुने गए अंतराल
से मिला था: और किसी लंबे-या-बराबर अंतराल से मिलने वाला अंतराल उसके
तिगुने $B \subseteq
3B_{i_l}$ में अंतर्विष्ट होता है।

\textbf{12.} $\{Mf > t\}$ विवृत है: प्रत्येक औसत $x \mapsto
\frac1{2r}\int_{x-r}^{x+r}\abs f$ \omterm{def:b3:topology:continuity}{संतत} है ($x$ में
प्रभावी अभिसरण), और \omterm{def:b3:topology:continuity}{संतत} फलनों का कोई उच्चतम अधःसंतत होता है। उसका
प्रत्येक $x$ $\int_{B_x}\abs f >
t\,\lambda(B_x)$ वाला $B_x =
\intoo{x-r_x}{x+r_x}$ रखता है। \omterm{def:b3:topology:compact}{संहत} $K \subseteq \{Mf > t\}$ के लिए: परिमित
संख्या के $B_x$ $K$ को आच्छादित कर लेते हैं, विटाली (प्रश्न 11)
$K \subseteq
\bigcup_l3B_l'$ वाले असंयुक्त $B_1', \dots, B_k'$ निकाल लेते हैं, अतः असंयुक्तता से
\[
\lambda(K) \leq 3\sum_l\lambda(B_l') <
\frac3t\sum_l\int_{B_l'}\abs f \leq \frac3t\,\norm f_1
\]
और आंतरिक नियमितता (\cref{thm:b3:measure:regularity}) से बात पूरी हो जाती है।

\textbf{13.} $x > 1$ के लिए: $r \in \intcc{x-1}x$ के साथ औसत $\frac{r - x + 1}{2r}$ है, जो $r$ में
बढ़ता है; और $r
\geq x$ के लिए वह $\frac1{2r}$ है, जो घटता है: अतः उच्चतम $\frac1{2x}$
है, जो $r = x$ पर प्राप्त होता है। इसलिए $M\mathbf 1_{\intcc01}
\notin L^1$। व्यापक रूप से, यदि
किसी परिबद्ध अंतराल $I \subseteq \intcc{-C}C$ पर $\int_I\abs f = c > 0$ हो, तो सभी $x$ के लिए $Mf(x) \geq
\frac{c}{2(\abs x + C)}$:
अतः जब तक लगभग सर्वत्र $f = 0$ न हो तब तक कभी \omterm{def:b3:lebesgue:l1}{समाकलनीय} नहीं।

\textbf{14.} $f_t = f\,\mathbf 1_{\abs f > t/2}$ के साथ: $M(f - f_t) \leq \frac t2$, अतः $\{Mf > t\} \subseteq \{Mf_t
> \frac t2\}$ और प्रश्न 12 $\lambda(Mf > t) \leq
\frac6t\int_{\abs f > t/2}\abs f$ देता
है। परत-केक (\cref{prop:b3:product:layercake}) और टोनेली:
\[
\norm{Mf}_p^p = p\int_0^\infty t^{p-1}\lambda(Mf > t)\dd t
\leq 6p\int\abs f\int_0^{2\abs f}t^{p-2}\,\dd t\,\dd\lambda
= \frac{6p\,2^{p-1}}{p-1}\,\norm f_p^p .
\]

\textbf{15.} $A_rf(x) = \frac1{2r}\int_{x-r}^{x+r}f$ लिखिए। \omterm{def:b3:topology:continuity}{संतत} $g$ के लिए: सर्वत्र $A_rg(x) \to g(x)$। $\varepsilon > 0$
दिया हो, तो \omterm{def:b3:topology:compact}{संहत} आलंब वाले \omterm{def:b3:topology:continuity}{संतत} $g$ और $\norm h_1 < \varepsilon$ (\cref{thm:b3:lp:density}) के साथ $f = g + h$
में विभाजित कीजिए; तब
\[
\limsup_{r\to0}\,\abs{A_rf(x) - f(x)} \leq Mh(x) +
\abs{h(x)},
\]
अतः $\Omega_\delta = \{\limsup_r\abs{A_rf - f} > \delta\}
\subseteq \{Mh > \tfrac\delta2\} \cup \{\abs h >
\tfrac\delta2\}$ का \omterm{def:b3:measure:measure}{माप} $\leq \frac{6\varepsilon}\delta
+ \frac{2\varepsilon}\delta$ है (प्रश्न 12; मार्कोव)। $\varepsilon$ स्वेच्छ: $\lambda(\Omega_\delta) = 0$;
और $\delta = \frac1k$ पर संघ: लगभग सर्वत्र $A_rf \to f$।

\textbf{16.} (क) प्रत्येक $q \in \Q$ के लिए $\abs{f - q}$ पर लगाया गया प्रश्न
15 लगभग सर्वत्र $A_r\abs{f - q}(x) \to \abs{f(x) - q}$ देता है; और इन पूर्ण-माप समुच्चयों के प्रतिच्छेद
पर $\abs{f(x) - q} < \eta$ वाला $q$ चुनने पर: प्रत्येक $\eta$ के लिए $\limsup_rA_r\abs{f - f(x)}(x) \leq 2\eta$: अर्थात्
लगभग प्रत्येक $x$ कोई लेबेग बिंदु है। (ख) किसी लेबेग बिंदु पर
\[
\Bigl|\frac{F(x + h) - F(x)}h - f(x)\Bigr|
= \Bigl|\frac1h\int_x^{x+h}\bigl(f - f(x)\bigr)\Bigr|
\leq 2\,A_{\abs h}\abs{f - f(x)}(x) \to 0 :
\]
लगभग सर्वत्र $F' = f$ --- अर्थात् $L^1$ फलनों के आद्यंतर वस्तुतः वापस
अवकलित हो जाते हैं; और सीढ़ी (\cref{pb:b3:measure:1}) केवल \emph{विलोम} दिशा का
प्रतिउदाहरण है।

\textbf{17.} प्रश्न 15 को $\mathbf 1_{A\cap[-n,n]}$ पर लगाइए और $n$ को बढ़ने दीजिए:
लगभग प्रत्येक $x \in A$ के लिए \omterm{ex:b3:lebesgue:gamma}{घनत्व} $\frac{\lambda(A\cap\intcc{x-r}{x+r})}{2r} \to 1$। श्टाइनहाउस: किसी \omterm{ex:b3:lebesgue:gamma}{घनत्व} बिंदु
के चारों ओर \omterm{ex:b3:lebesgue:gamma}{घनत्व} $> \frac34$ वाला $r$ लीजिए; $\abs t < \frac r2$ के लिए $A$ और
$A + t$ प्रत्येक लंबाई $\leq \frac52r$ के किसी अंतराल का $\frac32r$ से अधिक भाग भर
देते हैं, अतः वे प्रतिच्छेद करते हैं: $\intoo{-r/2}{r/2} \subseteq A - A$।

\textbf{18.} मान लीजिए $G(x) = x^{\alpha+1-p}F(x)^p$: $G$ $0$ पर लुप्त होता है ($0$
के निकट $F$ लुप्त होता है) और $\infty$ पर भी ($F$ परिबद्ध, $\alpha + 1 - p < 0$),
अतः $\int_0^\infty G' = 0$ जिसमें
\[
G' = (\alpha + 1 - p)\,x^{\alpha-p}F^p +
p\,x^{\alpha+1-p}F^{p-1}f
= (\alpha + 1 - p)\Bigl(\frac Fx\Bigr)^px^\alpha +
p\Bigl(\frac Fx\Bigr)^{p-1}f\,x^\alpha .
\]
इसलिए $\int(F/x)^px^\alpha = \frac{p}{p-1-\alpha}
\int(F/x)^{p-1}f\,x^\alpha$; और \omterm{def:b3:measure:measure}{माप} $x^\alpha\dd x$ के लिए होल्डर (घातांक $\frac p{p-1}$ और $p$)
प्रश्न 3 की भाँति बात पूरी कर देता है। सीमांत $\alpha = p - 1$: $f(t) =
\frac1t\mathbf 1_{\intcc1A}$ के साथ
दायाँ पक्ष $\ln A$ है जबकि बाएँ में $\int_1^A\frac{(\ln x)^p}x\dd x =
\frac{(\ln A)^{p+1}}{p+1}$ है: $A \to
\infty$ पर कोई अचर नहीं
बचता।

\textbf{19.} $\{0 < t < x\}$ पर टोनेली:
\[
\langle Hf, g\rangle =
\int_0^\infty\frac{g(x)}x\int_0^xf(t)\,\dd t\,\dd x
= \int_0^\infty f(t)\int_t^\infty\frac{g(x)}x\,\dd x\,\dd t
= \langle f, H^*g\rangle .
\]
द्वैतता: $\norm{H^*f}_p = \sup\{\langle f, Hg\rangle : g
\geq 0, \norm g_q \leq 1\} \leq \norm f_p\cdot
\sup\norm{Hg}_q \leq \frac q{q-1}\norm f_p = p\,\norm f_p$, जिसमें हार्डी को $L^q$ में लगाया गया है और जिसका
अचर $\frac q{q-1}$ $p$ के बराबर है।

\textbf{20.} (अकिंचन) विकर्ण के अनुदिश विभाजित कीजिए। $\{y \leq
x\}$ पर:
\[
\iint_{y\leq x}\frac{f(x)g(y)}{x}\,\dd y\,\dd x
= \int f(x)\,(Hg)(x)\,\dd x \leq \norm f_p\,\norm{Hg}_q
\leq p\,\norm f_p\norm g_q,
\]
क्योंकि $L^q$ में हार्डी अचर $\frac q{q-1} =
p$ ढोते हैं। सममित रूप से $\iint_{y>x} = \int g\,(Hf) \leq
\norm g_q\norm{Hf}_p \leq q\,\norm f_p\norm g_q$
($\frac
p{p-1} = q$)। कुल: $(p + q)\,\norm f_p\norm g_q$।

\textbf{21.} $a_n = n^{-1/p}$ के लिए $n \leq N$: दायाँ पक्ष $\bigl(\frac p{p-1}\bigr)^p\sum_{n\leq N}\frac1n =
\bigl(\frac p{p-1}\bigr)^p\ln N + O(1)$ है। बाईं ओर, $n
\leq N$
के लिए: $a_1 + \dots + a_n \geq \int_1^{n+1}t^{-1/p}\dd t =
\frac{p}{p-1}\bigl((n+1)^{1-1/p} - 1\bigr)$, अतः $n$-वाँ चेज़ारो माध्य किसी भी परास $n \geq n_0(\eta)$ में
$n$ के लिए एकसमान रूप से $\geq \frac p{p-1}n^{-1/p}(1 -
o(1))$ है; $p$ घात लेकर योग करने पर बायाँ
पक्ष $\geq
\bigl(\frac p{p-1}\bigr)^p(1 - \eta)\ln N + O_\eta(1)$ है। भाग देकर पहले $N \to \infty$ और फिर $\eta \to 0$ लेने पर: $\bigl(\frac p{p-1}\bigr)^p$ से
छोटा कोई अचर काम नहीं कर सकता।

\textbf{22.} $L^p$ पर $H$ परिबद्ध: संचयी औसत $p$-मानकों को नहीं
फुलाते (अचर $\frac p{p-1}$); $\ell^p$ पर चेज़ारो: वही, विविक्त रूप में; $L^p$
पर $M$ परिबद्ध: सबसे अच्छा स्थानीय औसत भी नियंत्रण में रहता है
(अचर $O(\frac1{p-1})$)। तीनों $p = 1$ पर विफल हो जाते हैं, और एक ही कारण से:
द्रव्यमान की किसी संकेंद्रित इकाई का औसत लेने से $\frac1x$ पुच्छ उत्पन्न
होती है (प्रश्न 10 और 13), और $\frac1x$ अनंत के निकट $L^1$ को छोड़कर
प्रत्येक $L^p$ में होता है। मृदुकरण द्रव्यमान को ठीक \omterm{def:b3:lebesgue:l1}{समाकलनीयता} के
प्रसंवादी सीमांत तक फैला देता है।

\textbf{23.} $b_n = a_n^{1/p}$ रखिए, अतः $\sum b_n^p = \sum
a_n$। प्रश्न 8
\[
\sum_{n\geq1}\Bigl(\frac{b_1 + \dots + b_n}n\Bigr)^{p}
\leq \Bigl(\frac p{p-1}\Bigr)^{p}\sum_{n\geq1}a_n,
\]
देता है, और समांतर माध्य--गुणोत्तर माध्य प्रत्येक योज्य पद को नीचे
से परिबद्ध कर देता है:
\[
\Bigl(\frac{b_1 + \dots + b_n}n\Bigr)^{p}
\geq \bigl(b_1\cdots b_n\bigr)^{p/n}
= \bigl(a_1\cdots a_n\bigr)^{1/n}.
\]
अतः \emph{प्रत्येक} $p > 1$ के लिए $\sum(a_1\cdots a_n)^{1/n} \leq \bigl(\frac
p{p-1}\bigr)^p\sum a_n$। $m =
p - 1$ के साथ
\[
\Bigl(\frac p{p-1}\Bigr)^{p} =
\Bigl(1 + \frac1m\Bigr)^{m+1},
\]
जो $m \to \infty$ पर घटकर $\eu$ तक जाता है ($\eu$ के लिए शास्त्रीय एकदिष्ट
ऊपरी अनुक्रम)। $p$ पर निम्नतम लेने से अचर $\eu$ के साथ कार्लमान
की असमिका मिल जाती है।

\textbf{24.} $a_n = \frac1n$ के लिए $(a_1\cdots a_n)^{1/n} =
(n!)^{-1/n}$। \cref{pb:b3:product:1} का परिबंधन $\sqrt{2\pi n}\,(n/\eu)^n \leq n! \leq \sqrt{2\pi
n}\,(n/\eu)^n\eu^{1/(12n)}$ देता है, अतः
\[
(n!)^{1/n} = \frac n\eu\,(2\pi n)^{1/(2n)}
\eu^{O(1/n^2)}
= \frac n\eu\Bigl(1 +
O\Bigl(\frac{\ln n}{n}\Bigr)\Bigr),
\]
क्योंकि $(2\pi n)^{1/(2n)} = \exp\bigl(\frac{\ln(2\pi
n)}{2n}\bigr)$। प्रतिलोमन करने पर $(n!)^{-1/n} = \frac\eu n(1 +
O(\frac{\ln n}n))$, और $n \leq N$ पर योग लेने पर:
\[
\sum_{n\leq N}(n!)^{-1/n}
= \eu\ln N + O(1),
\qquad
\eu\sum_{n\leq N}\frac1n = \eu\ln N + O(1)
\]
(त्रुटि श्रेणी $\sum\frac{\ln n}{n^2}$ अभिसरित होती है)। अचर $c$ वाली कोई कार्लमान
असमिका $\eu\ln N
+ O(1) \leq c\,(\ln N + O(1))$ अनिवार्य कर देती, अतः $\ln N$ से भाग देने पर $c \geq \eu$।
अनुकूलक अनुक्रम पंक्तिबद्ध हो जाते हैं: हार्डी के अचर तक $n^{-1/p}$ से
पहुँचा जाता है (प्रश्न 21), और कार्लमान के अचर तक $n^{-1}$ से --- और
प्रत्येक दशा में वह प्रसंवादी-प्रकार का अनुक्रम जो जिस समष्टि का
औसत लिया जा रहा है उसके ठीक बाहर बैठा हो।

\textbf{25.} $f = \mathbf 1_{\intcc01}$ के लिए: $\intoc01$ पर $Hf(x) = 1$ और $x > 1$ के लिए $Hf(x) = \frac1x$,
अतः $\norm{Hf}_2^2 = 1 + \int_1^\infty x^{-2}\dd x = 2$ और अनुपात $\sqrt2 \approx 1.414$ है, जो तीक्ष्ण परिबंध का लगभग $71\%$ है।
$f_A$ ($p = 2$) के लिए: $\intcc1A$ पर $F(x) = 2(\sqrt x - 1)$, अतः उस परास पर $Hf_A = 2x^{-1/2} - 2x^{-1}$ और, $x > A$
के लिए, $Hf_A(x) = 2(\sqrt A - 1)/x$। वर्ग करके समाकलन करने पर
\begin{align*}
\int_1^A\bigl(2x^{-1/2} - 2x^{-1}\bigr)^2\dd x
&= 4\ln A - 12 + \frac{16}{\sqrt A} - \frac4A,\\
\int_A^{\infty}\frac{4(\sqrt A - 1)^2}{x^2}\,\dd x
&= 4 - \frac{8}{\sqrt A} + \frac4A,
\end{align*}
जिससे $\norm{Hf_A}_2^2 = 4\ln A - 8 + 8A^{-1/2}$; और $\norm{f_A}_2^2 = \ln A$ से भाग देने पर कहा गया समीकरण मिल जाता है।
$A
= \eu^{10}$ पर: $4 - \frac{8(1 - \eu^{-5})}{10} = 3.2054$, अतः अनुपात $\sqrt{3.2054} \approx 1.790 < 2$ है। दोष $4 - \norm{Hf_A}_2^2/\norm{f_A}_2^2 \sim 8/\ln A$ केवल लघुगणकीय रूप
से क्षय होता है: अनुपात $1.99$ तक पहुँचने के लिए $\ln A \approx 200$ चाहिए होता,
अर्थात् $A \approx 10^{87}$। तीक्ष्ण अचर प्रमेय हैं, प्रयोग नहीं।
\end{solution}
